% !TEX root = ../main.tex
%\vspace{-3mm}
%\begin{abstract}
%  Diffusion Language Models (DLMs) have recently achieved strong results in text generation. However, their multi-step sampling leads to slow inference, limiting practical use. To address this, we extend Inverse Distillation, a technique originally developed to accelerate continuous diffusion models, to the discrete setting.
%  Nonetheless, this extension introduces both theoretical and practical challenges.
%From a theoretical perspective, the inverse distillation objective lacks uniqueness guarantees,
%which may lead to suboptimal solutions. From a practical standpoint, backpropagation in the discrete space is
%non-trivial and often unstable. To overcome these challenges, we first provide a theoretical result demonstrating
%that our inverse formulation admits a unique solution, thereby ensuring valid optimization. We then introduce gradient-stable relaxations to support effective training.
%  As a result, experiments on multiple DLMs show that our method, \textit{Inverse-distilled Diffusion Language Models (IDLM)}, reduces the number of inference steps by $4 \times$-$64 \times$, while preserving the teacher model’s entropy and generative perplexity.
%  \vspace{-1mm}
%\end{abstract}
%\begin{abstract}
%Diffusion language models (DLMs) enable competitive text generation, but their reverse diffusion typically requires hundreds to thousands of steps, making inference-time guidance, constraint satisfaction, and preference-oriented steering prohibitively expensive. We propose \emph{Inverse-distilled Diffusion Language Models} (IDLM), a post-training framework that distills a pretrained DLM teacher into a few-step generator by extending inverse distillation to discrete diffusion objectives.
%The discrete extension raises two obstacles: (i) the inverse objective may admit non-unique minimizers, and (ii) gradients through discrete sampling are unstable. We address (i) by proving that for common DLM objectives (SEDD, MDLM, and Duo) the IDLM loss lower-bounds \(D_{\mathrm{KL}}(p_\theta \,\|\, p^*)\) and has a unique global optimum at the data distribution. For (ii) we introduce gradient-stable relaxations via a simplex-parameterized generator and principled gradient paths through the discrete forward process.
%Across multiple pretrained DLMs on OpenWebText, IDLM achieves \(4\times\)–\(256\times\) fewer sampling steps while preserving teacher generative perplexity and sequence entropy, substantially improving the practicality of diffusion LMs in control-heavy, low-latency regimes.
%\end{abstract}
%\vspace{-3mm}
%\begin{abstract}
%Diffusion Language Models (DLMs) generate text via iterative reverse diffusion, but the resulting inference latency
%limits practical use and makes inference-time methods such as guidance expensive.
%We propose \emph{Inverse-distilled Diffusion Language Models (IDLM)}, a post-training framework that distills
%a pretrained DLM into a few-step generator by extending inverse distillation to discrete token spaces.
%IDLM uses a bilevel objective: an auxiliary \emph{fake} diffusion model is trained on samples from the student
%generator using the same diffusion loss as the teacher, and the student is updated to maximize the teacher-fake loss gap
%on its own samples. Two obstacles arise in discrete settings: the inverse objective may be non-identifiable,
%and gradients through discrete sampling can be unstable.
%We address identifiability by proving a uniqueness guarantee for IDLM under SEDD, MDLM, and Duo objectives.
%For stable optimization, we introduce gradient-friendly relaxations: the student outputs simplex-valued token distributions (avoiding brittle hard relaxations), and the diffusion losses are reformulated to remain computable and differentiable under this relaxation, enabling end-to-end training.
%As a result, experiments on multiple DLMs show that our method reduces the number of inference steps by $4 \times$-$64 \times$,
%while preserving the teacher model’s entropy and generative perplexity.
%\vspace{-2mm}
%\end{abstract}

\vspace{-3mm}
\begin{abstract}
Diffusion Language Models (DLMs) generate text via iterative reverse diffusion, but the resulting inference latency
limits practical use and makes inference-time methods such as guidance expensive.
We propose \emph{Inverse-distilled Diffusion Language Models (IDLM)}, a post-training framework that distills
a pretrained DLM into a few-step generator by extending inverse distillation to discrete token spaces.
IDLM optimizes a bilevel objective: a \emph{fake} diffusion model is trained on student samples with the teacher’s diffusion loss, and the student is updated to maximize the teacher--fake loss gap on its own samples.
In discrete settings, we (i) establish identifiability by proving a uniqueness guarantee under SEDD, MDLM, and Duo objectives, and (ii) stabilize training with simplex-valued token outputs and differentiable reformulations of the diffusion losses.
As a result, experiments on multiple DLMs show that our method reduces the number of inference steps by $4 \times$-$64 \times$,
while preserving the teacher model’s entropy and generative perplexity.
\vspace{-1mm}
\end{abstract}